\documentclass[11pt,a4paper]{article}

% This file is standalone for convenient review.  The content beginning at
% \section{Experimental evaluation} can be inserted directly after Section 3
% of problem_statement.pdf.
\usepackage[T1]{fontenc}
\usepackage[english]{babel}
\usepackage{lmodern}

\usepackage[a4paper,margin=25mm]{geometry}
\usepackage{microtype}
\usepackage{amsmath,amssymb}
\usepackage{booktabs}
\usepackage{tabularx}
\usepackage{array}
\usepackage{multirow}
\usepackage{graphicx}
\usepackage{placeins}
\usepackage{xcolor}
\usepackage{colortbl}
\usepackage{fancyhdr}
\usepackage{titlesec}
\usepackage{xurl}
\usepackage{hyperref}

\definecolor{navy}{HTML}{183153}
\definecolor{blue}{HTML}{2B6CB0}
\definecolor{lightblue}{HTML}{EDF5FC}
\definecolor{lightgray}{HTML}{F3F4F6}
\definecolor{midgray}{HTML}{697386}
\definecolor{green}{HTML}{287A55}
\definecolor{red}{HTML}{A43D3D}

\hypersetup{
  colorlinks=true,
  linkcolor=navy,
  urlcolor=blue,
  citecolor=green,
  pdfauthor={GMM research project},
  pdftitle={Experimental evaluation of localized moments for Gaussian mixtures},
  pdfsubject={Clean, contaminated, and heavy-tailed mixture estimation}
}

\pagestyle{plain}

\titleformat{\section}
  {\Large\bfseries}{\thesection}{1em}{}
\titleformat{\subsection}
  {\large\bfseries}{\thesubsection}{1em}{}
\titleformat{\subsubsection}
  {\normalsize\bfseries}{\thesubsubsection}{1em}{}
\titlespacing*{\section}{0pt}{1.30em}{0.55em}
\titlespacing*{\subsection}{0pt}{0.95em}{0.35em}
\titlespacing*{\subsubsection}{0pt}{0.70em}{0.25em}

\renewcommand{\arraystretch}{1.12}
\setlength{\tabcolsep}{4.0pt}
\setlength{\parindent}{0pt}
\setlength{\parskip}{0.35em}
\emergencystretch=2em

\newcolumntype{Y}{>{\raggedright\arraybackslash}X}
\newcolumntype{R}{>{\raggedleft\arraybackslash}X}
\newcommand{\best}[1]{\textbf{#1}}
\newcommand{\bad}[1]{\textcolor{red}{#1}}
\newcommand{\E}{\mathbb{E}}
\newcommand{\KL}{\mathrm{KL}}
\newcommand{\ARI}{\mathrm{ARI}}
\newcommand{\Mzero}{M_0}
\newcommand{\Mone}{M_1}
\newcommand{\Mtwo}{M_2}
\newcommand{\Rtwo}{R_2}
\newcommand{\EM}{\mathrm{EM}}

\begin{document}
\setcounter{section}{3}
\setcounter{page}{12}

\section{Experimental evaluation}
\label{sec:experiments}

We evaluate whether localized moment matching can recover a spherical
Gaussian mixture under correct specification and under several forms of
model misspecification.  The experiments are organized around three
questions.  First, which moment families are informative when the Gaussian
model is correct?  Second, does the local criterion remain useful when a
small fraction of observations does not belong to any component?  Third,
which parts of the result are specific to uniform background noise, and which
extend to global heavy tails?  We also isolate the effects of radial moments
and alternative variance-update schemes.  Unless stated otherwise, all
reported errors are computed after optimal permutation of the components.

\subsection{Experimental protocol}
\label{sec:exp_protocol}

\subsubsection{Data-generating processes}

The core problem has $K=3$ balanced spherical components in dimension
$d=10$ and $n=200$ observations.  Component means are generated with a
minimum pairwise Mahalanobis separation of $2$, producing deliberately
overlapping clusters.  The default ratio between the largest and smallest
component variance is $3$.  Table~\ref{tab:datasets} summarizes the data
regimes.  Each dataset seed defines the mixture parameters and observations;
model seeds change only randomized directions and optimization details.

\begin{table}[ht]
\centering
\small
\caption{Data regimes used in the experiments.  The target distribution for
evaluation is always the uncontaminated Gaussian mixture.}
\label{tab:datasets}
\begin{tabularx}{\linewidth}{@{}l l Y Y@{}}
\toprule
Regime & Grid & Construction & Purpose \\
\midrule
Clean Gaussian
  & variance ratio $1,1.2,3,6$
  & Samples from the target spherical GMM.
  & Calibration under correct specification. \\
Uniform background
  & $\varepsilon=0,2,5,10\%$
  & Replace an $\varepsilon$ fraction by a uniform draw from the bounding box
    of the mixture, enlarged by padding $4$.
  & Localized contamination and negative-space information. \\
Student-$t$
  & $\nu=10,5,3$
  & Replace every Gaussian component by a moment-matched spherical
    Student-$t_\nu$ component; labels and covariance targets are retained.
  & Symmetric global heavy tails. \\
Directional Gamma
  & 12 settings
  & Skewed component-wise perturbations with random, outward, or inward
    directions; $n\in\{200,500\}$ and two separations.
  & Asymmetric misspecification and a negative control. \\
Dimension scaling
  & $d=2,10,50,100,500$
  & Clean and $5\%$ uniform data with
    $n(d)=\max(200,20d)$ and a fixed test-function budget.
  & Statistical and computational scaling. \\
\bottomrule
\end{tabularx}
\end{table}

The clean calibration uses 20--30 independent datasets and two model seeds.
The main uniform-contamination and Student-$t$ experiments use 60 independent
datasets and three model seeds.  Smaller mechanistic and optimizer ablations
use 20--40 datasets and two or three model seeds.  The dimension screen uses
20 datasets and one model seed in each cell.  Exact sample sizes are given
with the corresponding tables.

\subsubsection{Moment models}

For each center $c$, bandwidth $s$, and unit direction $u$, the models combine
the normalized zeroth-, first-, and directional second-order tests introduced
in Section~3,
\[
  \Mzero=K_s(x-c),\qquad
  \Mone=\frac{\langle x-c,u\rangle}{s}K_s(x-c),\qquad
  \Mtwo=\frac{\langle x-c,u\rangle^2}{s^2}K_s(x-c).
\]
We use all observations as centers and leave the center itself out when
forming its empirical response and local neighbor count.  This leave-one-out
construction removes the self-interaction bias caused by data-centered test
functions.  The default bandwidths contain approximately 17 and 133
neighbors.  Ten random directions are generated for both $\Mone$ and
$\Mtwo$.  All configurations build the same $1+10+10$ test stream per
center and bandwidth; a family is disabled by assigning zero loss weight.
Thus, differences between moment configurations do not arise from different
random features or computational budgets.

The notation $\Mtwo$ below means ten directional second moments at both
bandwidths.  For example, $\Mtwo+0.25\Mone$ assigns unit weight to the
second-order block and weight $0.25$ to the first-order block.  The radial
second moment is
\[
 \Rtwo(x;c,s)=\frac{\|x-c\|^2}{d s^2}K_s(x-c),
\]
with one radial test replacing the directional second-order block.  Its
compensated version subtracts the multiple of $\Mzero$ derived in Section~3
so that the test integrates to zero over $\mathbb R^d$.  Division by $d$ and
$s^2$ puts radial and directional responses on comparable scales.

\subsubsection{Baselines, optimization, and metrics}

The baselines are spherical EM, initialized by KMeans and fitted with
$\texttt{reg\_covar}\in\{10^{-6},10^{-2}\}$, and KMeans itself.  The moment
estimator uses the same KMeans initialization.  Mixture weights are solved on
the simplex.  Geometry is then updated component by component: with all other
components fixed, the component mean and log-variance are optimized jointly
by a bounded trust-region least-squares step.  Mean displacement is restricted
relative to the current component scale by a step bound of $0.75$; a local step is accepted only if it
does not increase the moment loss.  Weight and geometry sweeps alternate for
at most 150 outer iterations, with relative tolerance $10^{-5}$ and patience
three.  Section~\ref{sec:optimizer_ablation} tests both stricter stopping and
two explicitly separated mean/variance updates.

We report forward $\KL(P^*\|P_{\hat\theta})$ from the clean target GMM to the
estimate, total variation error of the mixture weights, mean error normalized
by the target scale, affine-invariant covariance error, and adjusted Rand
index (ARI).  Under contamination, parameter errors and KL refer to the clean
target, while ARI is evaluated only on genuine observations.  Lower is better
for all errors and higher is better for ARI.

In main tables, an entry is the mean with the standard deviation in
parentheses.  We first average model seeds within each dataset, then summarize
over independent dataset seeds; the experimental unit is therefore the
dataset, not an optimizer restart.  Calibration tables for which only the
original aggregate output was retained summarize all completed fits.  Paired
differences use dataset-level percentile bootstrap 95\% confidence intervals
(CI).  No multiplicity correction is applied, so fine-grained weight searches
are interpreted as exploratory rather than as confirmatory hypothesis tests.

\subsection{Calibration under the Gaussian model}
\label{sec:clean_calibration}

The leave-one-out correction is consequential in the nearly homoscedastic
regime: for the wide-bandwidth $\Mzero+3\Mone$ model it reduces mean KL from
$0.100$ to $0.087$ at variance ratio $1$, and from $0.101$ to $0.088$ at ratio
$1.2$.  We therefore use leave-one-out responses in every subsequent
experiment.

Table~\ref{tab:clean} compares the best calibrated moment model with EM and
KMeans.  Local moments are best when component variances are nearly equal.
At the default variance ratio $3$, $\Mtwo+0.25\Mone$ is essentially tied with
EM: it has slightly larger KL and covariance error, but smaller weight error
and slightly larger ARI.  At ratio $6$, EM is clearly preferable.  Thus, the
clean-data result is not a uniform dominance claim; it identifies where the
localized criterion pays little or no efficiency penalty.

\begin{table}[ht]
\centering
\scriptsize
\caption{Clean Gaussian data.  Ratios $1$, $1.2$, and $6$ use mean (SD) over
20 dataset seeds and two model seeds; ratio $3$ uses 60 datasets after
averaging three model seeds.  Bold denotes the better estimator in each pair.}
\label{tab:clean}
\resizebox{\linewidth}{!}{%
\begin{tabular}{@{}c l c c c c c@{}}
\toprule
Variance ratio & Estimator & KL $\downarrow$ & Weight TV $\downarrow$
  & Mean error $\downarrow$ & Covariance $\downarrow$ & ARI $\uparrow$ \\
\midrule
\multirow{2}{*}{$1.0$}
 & $\Mzero+3\Mone$, $s(133)$ & \best{.0867 (.0281)} & \best{.0536 (.0311)}
 & \best{.1846 (.0398)} & \best{.0511 (.0274)} & \best{.3263 (.0726)} \\
 & EM, $10^{-2}$ & .1076 (.0382) & .1329 (.0783) & .2245 (.0657)
 & .1122 (.0640) & .2894 (.0822) \\
\addlinespace
\multirow{2}{*}{$1.2$}
 & $\Mzero+3\Mone$, $s(133)$ & \best{.0883 (.0244)} & \best{.0491 (.0232)}
 & \best{.1812 (.0383)} & \best{.0589 (.0244)} & \best{.3548 (.0563)} \\
 & EM, $10^{-2}$ & .1079 (.0355) & .1356 (.0756) & .2166 (.0590)
 & .1204 (.0633) & .2985 (.0763) \\
\addlinespace
\multirow{2}{*}{$3.0$}
 & $\Mtwo+0.25\Mone$ & .1049 (.0269) & \best{.0590 (.0243)}
 & .1590 (.0236) & .0842 (.0383) & \best{.5795 (.0604)} \\
 & EM, $10^{-2}$ & \best{.0989 (.0289)} & .0746 (.0407)
 & \best{.1564 (.0269)} & \best{.0825 (.0470)} & .5729 (.0643) \\
\addlinespace
\multirow{2}{*}{$6.0$}
 & $\Mone+\Mtwo$, wide & .1182 (.0378) & .0632 (.0278)
 & .1530 (.0258) & .0991 (.0511) & .7402 (.0689) \\
 & EM, $10^{-2}$ & \best{.0901 (.0228)} & \best{.0487 (.0258)}
 & \best{.1348 (.0201)} & \best{.0667 (.0267)} & \best{.7698 (.0527)} \\
\bottomrule
\end{tabular}}
\end{table}

Two bandwidths, with approximately 17 and 133 neighbors, consistently
outperform either a single bandwidth or a denser three-bandwidth grid
$[17,67,133]$.  Directional coverage saturates quickly: increasing the number
of second-order directions from 1 to 3 to 10 changes KL from $0.1133$ to
$0.1075$ to $0.1058$, whereas 20 directions give $0.1057$ with no ARI gain.
We consequently fix ten directions and two bandwidths rather than tune them
separately in later experiments.

\subsection{Uniform background contamination}
\label{sec:uniform}

Uniform replacement noise is the regime in which localized second moments
are most useful.  Table~\ref{tab:uniform_main} reports the best configuration
from a small first-order weight grid at each contamination level.  On clean
data, EM retains a small KL advantage.  From $2\%$ to $10\%$ contamination,
the moment estimator has lower KL and higher ARI than both baselines.  At
$5\%$, for example, KL falls from $0.278$ for EM to $0.145$, while ARI rises
from $0.385$ to $0.534$.  The full set of parameter metrics at this central
noise level is shown in Table~\ref{tab:uniform_full}.

\begin{table}[ht]
\centering
\scriptsize
\caption{Uniform contamination: mean (SD) over 60 independent datasets after
averaging three model seeds.  The moment configuration is selected from the
predefined first-order weight grid.}
\label{tab:uniform_main}
\begin{tabular}{@{}c l c c c@{}}
\toprule
$\varepsilon$ & Estimator & KL $\downarrow$ & ARI $\uparrow$
  & Selected moment weights \\
\midrule
\multirow{3}{*}{$0\%$}
 & Local moments & .1049 (.0269) & \best{.5795 (.0604)} & $\Mtwo+0.25\Mone$ \\
 & EM, $10^{-2}$ & \best{.0989 (.0289)} & .5729 (.0643) & --- \\
 & KMeans & .2290 (.0561) & .4865 (.0824) & --- \\
\addlinespace
\multirow{3}{*}{$2\%$}
 & Local moments & \best{.1132 (.0303)} & \best{.5667 (.0605)} & $\Mtwo+0.125\Mone$ \\
 & EM, $10^{-2}$ & .1919 (.0450) & .4699 (.0981) & --- \\
 & KMeans & .2874 (.0640) & .4459 (.0795) & --- \\
\addlinespace
\multirow{3}{*}{$5\%$}
 & Local moments & \best{.1449 (.0384)} & \best{.5339 (.0800)} & $\Mtwo$ \\
 & EM, $10^{-2}$ & .2775 (.0381) & .3845 (.0649) & --- \\
 & KMeans & .4407 (.0753) & .3871 (.0880) & --- \\
\addlinespace
\multirow{3}{*}{$10\%$}
 & Local moments & \best{.3208 (.0656)} & \best{.4023 (.0978)} & $\Mtwo$ \\
 & EM, $10^{-2}$ & .3387 (.0501) & .3688 (.0779) & --- \\
 & KMeans & .7712 (.1151) & .2888 (.0912) & --- \\
\bottomrule
\end{tabular}
\end{table}

\begin{table}[ht]
\centering
\scriptsize
\caption{Uniform contamination at $\varepsilon=5\%$: full parameter recovery
metrics, mean (SD) over 60 datasets.}
\label{tab:uniform_full}
\begin{tabular}{@{}l c c c c c@{}}
\toprule
Estimator & KL $\downarrow$ & Weight TV $\downarrow$ & Mean error $\downarrow$
 & Covariance $\downarrow$ & ARI $\uparrow$ \\
\midrule
$\Mtwo$ & \best{.1449 (.0384)} & \best{.0873 (.0442)}
 & \best{.2091 (.2205)} & \best{.1889 (.0880)} & \best{.5339 (.0800)} \\
EM, $10^{-2}$ & .2775 (.0381) & .2917 (.0309)
 & .5019 (.1262) & 1.0959 (.2165) & .3845 (.0649) \\
KMeans & .4407 (.0753) & .1314 (.0479)
 & .3108 (.1606) & .4682 (.0876) & .3871 (.0880) \\
\bottomrule
\end{tabular}
\end{table}

The paired comparison in Table~\ref{tab:uniform_ci} confirms that the result
is not driven by a few easy datasets.  At every nonzero contamination level,
the 95\% CIs for both KL and ARI exclude zero.  At $10\%$, the KL improvement
has become small, but the clustering improvement remains detectable.

\begin{table}[ht]
\centering
\small
\caption{Dataset-paired difference $\Mtwo-\EM_{10^{-2}}$ with bootstrap 95\%
CI.  Negative KL and positive ARI favor localized moments.}
\label{tab:uniform_ci}
\begin{tabular}{@{}c c c@{}}
\toprule
$\varepsilon$ & $\Delta$KL [95\% CI] & $\Delta$ARI [95\% CI] \\
\midrule
$0\%$  & $+.0107$ $[+.0041,+.0170]$ & $+.0003$ $[-.0098,+.0104]$ \\
$2\%$  & \best{$-.0764$ $[-.0898,-.0630]$} & \best{$+.0959$ $[+.0728,+.1189]$} \\
$5\%$  & \best{$-.1326$ $[-.1437,-.1213]$} & \best{$+.1494$ $[+.1277,+.1715]$} \\
$10\%$ & \best{$-.0179$ $[-.0345,-.0015]$} & \best{$+.0334$ $[+.0098,+.0575]$} \\
\bottomrule
\end{tabular}
\end{table}

\subsubsection{Moment-order ablation}

The optimal first-order weight decreases monotonically with contamination:
$0.25$ on clean data, approximately $0.125$ at $2\%$, and zero at $5\%$ and
$10\%$.  On clean data, adding $0.25\Mone$ to $\Mtwo$ changes KL by
$-.0047$ (95\% CI $[-.0070,-.0029]$) and ARI by $+.0063$
($[+.0020,+.0108]$).  At $5\%$, even weight $0.03125$ changes KL by $+.0038$
and ARI by $-.0038$, with both CIs excluding zero; at $10\%$, the changes are
$+.0160$ and $-.0066$.  First moments encode local asymmetry and therefore
help locate a correctly specified Gaussian component, but they also react
directly to background observations.  The observed weight path is consistent
with this bias--variance trade-off.

Zeroth-order moments are more damaging.  Adding weight $1/16$ to $\Mzero$
gives KL/ARI $0.1430/0.5275$ at $2\%$, compared with $0.1155/0.5657$ for pure
$\Mtwo$; the paired differences, $+.0275$ and $-.0382$, have CIs excluding
zero.  At $5\%$ and $10\%$ the same model reaches KL $0.2861$ and $0.4258$.
The zeroth response measures local mass itself, so diffuse outliers create a
systematic model--data discrepancy everywhere.  Local second moments are not
immune to that discrepancy, but the factor $\langle x-c,u\rangle^2K_s(x-c)$
suppresses observations both at the center and far away, making the useful
signal more shell-like.

\subsubsection{Are outliers useful as test centers?}

One possible explanation for robustness is that outliers supply centers in
low-density regions and explicitly constrain the fitted mixture there.  We
test this at $5\%$ contamination by retaining all observations in the moment
averages while changing only the centers used to place the tests.
Table~\ref{tab:center_control} shows the opposite effect: replacing observed
centers by oracle clean centers yields a small improvement.  With bandwidths
held fixed, the paired KL change is $-.00511$ (95\% CI
$[-.00660,-.00389]$), covariance error changes by $-.01086$, and ARI by
$+.00446$; KL improves on 39 of 40 datasets.  Reselecting bandwidths from the
clean centers adds no detectable improvement ($\Delta$KL $=-.00063$, CI
includes zero).  Robustness therefore comes from the response and objective,
not from deliberately probing negative space with outlier centers.

\begin{table}[ht]
\centering
\small
\caption{Center-placement control at $5\%$ uniform contamination; means over
40 datasets and three model seeds.  All observations remain in the empirical
moment averages.}
\label{tab:center_control}
\begin{tabular}{@{}l c c c c c@{}}
\toprule
Centers / bandwidths & KL & Weight TV & Mean error & Covariance & ARI \\
\midrule
Observed / observed & .1496 & .0883 & .1871 & .1696 & .5528 \\
Oracle clean / fixed & \best{.1445} & \best{.0843} & \best{.1836} & .1587 & .5573 \\
Oracle clean / reselected & \best{.1438} & .0850 & .1860 & \best{.1578} & \best{.5577} \\
\bottomrule
\end{tabular}
\end{table}

\subsection{Scaling with dimension}
\label{sec:dimension_scaling}

We next vary $d\in\{2,10,50,100,500\}$ on clean data and at $5\%$ uniform
contamination.  To avoid the statistically ill-posed regime $n<d$, sample
size follows $n(d)=\max(200,20d)$, preserving the order of
$\sqrt{d/n_k}$ from the $d=10,n=200$ reference problem.  This is deliberately
a fixed-compute screen: every moment fit samples 100 data centers, uses ten
directions and the same two neighbor fractions $17/199$ and $133/199$.
We compare the three previously selected configurations $\Mtwo$,
$\Mtwo+0.03125\Mone$, and $\Mtwo+0.25\Mone$; Tables~\ref{tab:dim_clean}
and~\ref{tab:dim_uniform} display the one with the lowest mean KL in each
cell.  This within-grid choice is exploratory.

\begin{table}[ht]
\centering
\scriptsize
\caption{Clean dimension scaling: mean (SD) over 20 datasets and one model
seed.  The moment weight shown minimizes mean KL among the three prespecified
configurations.  KL is not divided by dimension, so comparisons are within a
row.}
\label{tab:dim_clean}
\resizebox{\linewidth}{!}{%
\begin{tabular}{@{}r r c c c c c c c@{}}
\toprule
$d$ & $n$ & Moment $w_1$ & Moment KL & EM KL & KMeans KL
 & Moment ARI & EM ARI & KMeans ARI \\
\midrule
2   & 200   & .25    & .0318 (.0121) & \best{.0287 (.0146)} & .0578 (.0150)
 & .4670 (.0688) & .4717 (.0751) & \best{.4800 (.0802)} \\
10  & 200   & .25    & .1100 (.0312) & \best{.0953 (.0257)} & .2154 (.0317)
 & .5650 (.0813) & \best{.5736 (.0680)} & .4873 (.0611) \\
50  & 1000  & .25    & .4138 (.0881) & \best{.0796 (.0097)} & .6365 (.0916)
 & .6612 (.0461) & \best{.8450 (.0147)} & .6817 (.0337) \\
100 & 2000  & .03125 & .7499 (.3017) & \best{.0818 (.0077)} & 1.179 (.0867)
 & .8375 (.0681) & \best{.9482 (.0094)} & .7901 (.0197) \\
500 & 10000 & .25    & 5.261 (4.218) & \best{1.040 (2.319)} & 6.692 (3.100)
 & .9057 (.1975) & \best{.9352 (.1583)} & .8187 (.1378) \\
\bottomrule
\end{tabular}}
\end{table}

\begin{table}[ht]
\centering
\scriptsize
\caption{Dimension scaling with $5\%$ uniform contamination.  Reporting and
selection are as in Table~\ref{tab:dim_clean}.}
\label{tab:dim_uniform}
\resizebox{\linewidth}{!}{%
\begin{tabular}{@{}r r c c c c c c c@{}}
\toprule
$d$ & $n$ & Moment $w_1$ & Moment KL & EM KL & KMeans KL
 & Moment ARI & EM ARI & KMeans ARI \\
\midrule
2   & 200   & .25    & \best{.0410 (.0117)} & .0853 (.0247) & .0853 (.0229)
 & .4553 (.1022) & .2700 (.1200) & \best{.4607 (.0864)} \\
10  & 200   & 0      & \best{.1524 (.0487)} & .2745 (.0464) & .4378 (.0782)
 & \best{.5375 (.1098)} & .3825 (.0897) & .3786 (.0796) \\
50  & 1000  & .03125 & 1.100 (.134) & \best{.6271 (.0169)} & 2.019 (.234)
 & .4306 (.0320) & \best{.4796 (.0251)} & .3521 (.0643) \\
100 & 2000  & 0      & 2.360 (1.052) & \best{1.192 (.0408)} & 4.098 (.418)
 & .4620 (.1100) & \best{.4919 (.0226)} & .3346 (.0285) \\
500 & 10000 & .03125 & 11.864 (6.538) & \best{6.805 (1.714)} & 19.523 (2.412)
 & \best{.6684 (.2133)} & .5678 (.0061) & .4967 (.0358) \\
\bottomrule
\end{tabular}}
\end{table}

The low-dimensional contamination advantage is reproduced at $d=2$ and
$d=10$, but reverses at $d=50$.  In a paired comparison of the fixed pure
$\Mtwo$ model with EM, the KL differences are $-.0441$
($[-.0546,-.0344]$), $-.1221$ ($[-.1398,-.1033]$), and $+.5050$
($[+.4238,+.5873]$) for $d=2,10,50$.  The corresponding ARI differences are
$+.1787$ ($[+.1155,+.2434]$), $+.1550$ ($[+.1176,+.1930]$), and $-.0679$
($[-.0860,-.0495]$).  At $d=100$ and $500$, pure-$\Mtwo$ KL remains
significantly worse, whereas its ARI difference from EM is inconclusive.
On clean data, the fixed $\Mtwo+0.25\Mone$ model is tied with EM at $d=2$ in
paired KL, is slightly worse at $d=10$, and is substantially worse thereafter.

The degradation is not a runtime failure: every fit returns successfully.
It is nevertheless statistically unstable.  Using the same scale/weight/mean
diagnostic as in Section~\ref{sec:heavy_tails}, the selected contaminated
moment model is flagged in $25\%$, $100\%$, and $30\%$ of runs at
$d=50,100,500$, compared with $0\%$, $0\%$, and $5\%$ for EM.  At $d=500$
clean data, all methods inherit a $15\%$ failure rate from their shared
KMeans initialization.  With a fixed 2,000 second-order conditions, the
moment objective is only mildly overdetermined relative to the 1,505 mixture
parameters at $d=500$, and the localized conditions remain correlated.  The
screen therefore shows that increasing $n$ proportionally to $d$ is not
sufficient: directional coverage, the number of centers, and initialization
must also scale.  The robustness claim in Section~\ref{sec:uniform} should be
restricted to low and moderate dimension under the present feature budget.

\subsection{Radial second moments}
\label{sec:radial}

Radial tests replace multiple directional projections by a single rotation-
invariant shell.  Table~\ref{tab:radial_clean} shows that the uncompensated
radial moment is effective on clean data.  With equal variances it improves
covariance recovery over a directional configuration, while the hybrid
$\Mone+4\Mtwo+4\Rtwo$ matches the fully directional model at variance ratio
$3$.  Compensation is consistently worse, despite its appealing zero-integral
property.  Although the identity cancels an ideal constant background at the
population level, the signed subtraction creates cancellation between two
noisy empirical quantities, is affected by the finite box boundary, and makes
the fitted objective more difficult to optimize.

\begin{table}[ht]
\centering
\scriptsize
\caption{Radial-moment ablation on clean data: mean (SD) over 20 datasets at
variance ratio $1$ and 30 at ratio $3$, with two model seeds.  ``Comp.''
denotes the zero-integral compensated radial test.}
\label{tab:radial_clean}
\begin{tabular}{@{}c l c c c@{}}
\toprule
Variance ratio & Moment model & KL $\downarrow$ & Covariance $\downarrow$ & ARI $\uparrow$ \\
\midrule
\multirow{3}{*}{$1$}
 & $3\Mone+8\Rtwo$ & \best{.0855 (.0257)} & \best{.0454 (.0206)} & .3256 (.0628) \\
 & $3\Mone+8\Rtwo$, comp. & .0916 (.0296) & .0522 (.0281) & \best{.3336 (.0738)} \\
 & $3\Mone+3\Mtwo$, wide & .1048 (.0284) & .0757 (.0214) & .3210 (.0607) \\
\addlinespace
\multirow{3}{*}{$3$}
 & $\Mone+4\Mtwo+4\Rtwo$ & \best{.1043 (.0277)} & .0896 (.0389) & .5788 (.0557) \\
 & same, comp. & .1213 (.0296) & .1245 (.0396) & .5502 (.0648) \\
 & $\Mone+8\Mtwo$ & .1058 (.0270) & \best{.0891 (.0397)} & \best{.5794 (.0585)} \\
\bottomrule
\end{tabular}
\end{table}

Under uniform contamination, however, the directional moment is preferable.
Table~\ref{tab:radial_contamination} compares a half-directional,
half-radial hybrid with pure $\Mtwo$.  The uncompensated hybrid has higher KL
and lower ARI at $2\%$ and $5\%$; at $10\%$ the KL difference is inconclusive,
but the ARI loss remains significant.  Compensation produces a large and
systematic deterioration: its analytical background cancellation does not
translate into better parameter recovery.  Radial moments are therefore a
useful compact clean-data feature, not the source of the principal robustness
result.

\begin{table}[ht]
\centering
\scriptsize
\caption{Radial moments under uniform contamination: mean (SD) over 60
datasets after averaging three model seeds.}
\label{tab:radial_contamination}
\begin{tabular}{@{}c l c c c@{}}
\toprule
$\varepsilon$ & Moment model & KL $\downarrow$ & Covariance $\downarrow$ & ARI $\uparrow$ \\
\midrule
\multirow{3}{*}{$2\%$}
 & $\Mtwo$ & \best{.1155 (.0307)} & \best{.1107 (.0456)} & \best{.5657 (.0611)} \\
 & $0.5\Mtwo+0.5\Rtwo$ & .1237 (.0321) & .1301 (.0487) & .5452 (.0663) \\
 & same, comp. & .1852 (.0476) & .2346 (.0606) & .4576 (.0819) \\
\addlinespace
\multirow{3}{*}{$5\%$}
 & $\Mtwo$ & \best{.1449 (.0384)} & \best{.1889 (.0880)} & \best{.5339 (.0800)} \\
 & $0.5\Mtwo+0.5\Rtwo$ & .1546 (.0397) & .2080 (.0969) & .5094 (.0853) \\
 & same, comp. & .3441 (.0510) & .4875 (.0875) & .3261 (.0970) \\
\addlinespace
\multirow{3}{*}{$10\%$}
 & $\Mtwo$ & \best{.3208 (.0656)} & \best{.5305 (.1883)} & \best{.4023 (.0978)} \\
 & $0.5\Mtwo+0.5\Rtwo$ & .3235 (.0642) & .5351 (.1918) & .3818 (.1020) \\
 & same, comp. & .5215 (.0659) & 1.2133 (1.4154) & .2126 (.0914) \\
\bottomrule
\end{tabular}
\end{table}

For the uncompensated hybrid relative to $\Mtwo$, paired KL differences are
$+.0082$ ($[+.0065,+.0099]$), $+.0097$ ($[+.0063,+.0136]$), and $+.0027$
($[-.0053,+.0109]$) at $2\%$, $5\%$, and $10\%$, respectively.  The
corresponding ARI differences are $-.0205$, $-.0245$, and $-.0205$, with all
three 95\% CIs below zero.

\subsection{Optimization ablations}
\label{sec:optimizer_ablation}

The nonconvex moment objective raises two distinct questions: whether the
default stopping rule leaves appreciable optimization error, and whether
joint mean/variance component updates are necessary.  We address them
separately.

\subsubsection{Stopping, polishing, and update granularity}

Table~\ref{tab:optimizer_standard} compares the default component-wise method
with a final joint geometry polish, a stricter iteration budget, and joint
optimization of all geometry from the initial solution.  Both polishing and
stricter stopping change the final objective by less than $0.1\%$ and have no
meaningful external-metric effect.  In contrast, optimizing all means and
variances jointly from the start is worse in both regimes and produces one
near-degenerate contaminated fit.  The result is consistent with poor basin
selection rather than premature termination.  A final polish is omitted from
the main experiments because its median relative loss gain is only
$3\times10^{-5}$ for $\Mtwo$ and $3\times10^{-4}$ for the mixed model.

\begin{table}[ht]
\centering
\scriptsize
\caption{Optimizer diagnostics: mean (SD) over 20 new datasets and two model
seeds.  Every row uses the same pure-$\Mtwo$ objective.}
\label{tab:optimizer_standard}
\begin{tabular}{@{}l c c c c c c@{}}
\toprule
Data & Geometry update & KL $\downarrow$ & Covariance $\downarrow$ & ARI $\uparrow$
 & Outer iters & Converged \\
\midrule
\multirow{4}{*}{Clean}
 & Component-wise & .1125 (.0345) & .0831 (.0375) & .5623 (.0822) & 79.5 & 1.00 \\
 & + joint polish & .1124 (.0345) & .0829 (.0377) & .5618 (.0823) & 79.5 & 1.00 \\
 & Strict budget & .1125 (.0345) & .0830 (.0376) & .5620 (.0832) & --- & --- \\
 & Joint from start & .1166 (.0366) & .0911 (.0392) & .5561 (.0870) & --- & --- \\
\addlinespace
\multirow{4}{*}{$5\%$ uniform}
 & Component-wise & \best{.1516 (.0481)} & \best{.1801 (.0729)} & .5420 (.1060) & 110.3 & .90 \\
 & + joint polish & \best{.1515 (.0481)} & \best{.1800 (.0729)} & .5428 (.1057) & 110.3 & .90 \\
 & Strict budget & \best{.1511 (.0479)} & \best{.1793 (.0722)} & \best{.5430 (.1056)} & --- & --- \\
 & Joint from start & .1612 (.0518) & .5198 (1.4625) & .5197 (.1060) & --- & --- \\
\bottomrule
\end{tabular}
\end{table}

The mixed $\Mtwo+0.25\Mone$ objective is harder to optimize: the stopping
criterion is reached in only $45.0\%$ of clean and $52.5\%$ of contaminated
fits, compared with $100\%$ and $90.0\%$ for pure $\Mtwo$.  Nevertheless,
extra iterations and polishing do not improve the external metrics.  Training
loss is therefore useful as an optimization diagnostic but not as a reliable
selector between different moment families.

The absolute loss is also easy to misinterpret.  At $5\%$ contamination the
mean final pure-$\Mtwo$ loss is $3.93\times10^{-6}$, but this is an average
squared residual over 4,000 second-order conditions.  Its root scale,
approximately $1.9\times10^{-3}$, should be compared with a typical empirical
small-bandwidth response of $3.1\times10^{-2}$.  The model has 35 free
parameters and the conditions are correlated; a small training loss is
therefore neither literal interpolation nor evidence of correct distribution
recovery.  In 120 controlled fits, final loss has correlations only $0.22$
with KL, $0.13$ with covariance error, and $0.03$ with ARI.  Leave-one-out
removes self-interaction bias but not optimism from evaluating the criterion
on the fitting sample.  A held-out moment loss remains an important future
model-selection diagnostic.

\subsubsection{Separating mean and variance updates}

We implemented two additional optimizers without changing the primary code
path.  The \emph{formula update} performs a mean-only component step and then
estimates its variance from the ratio of localized radial zeroth- and
second-order responses.  Before taking the ratio, it subtracts the modeled
contributions of every other component, as required by the correction in
Equation~(26); the resulting ratio is clipped only to keep the closed-form
variance finite.  Weights are re-estimated after a full component sweep.  The
\emph{joint-variance block} performs a complete mean sweep with all variances
fixed and then optimizes all log-variances jointly by least squares.

Table~\ref{tab:variance_optimizers} shows a sharp distinction.  The formula
update is competitive in the narrow homoscedastic setting, but fails for
variance ratio $3$ and under contamination: roughly 59--64\% of its outer
sweeps increase the objective, even though invalid or clipped ratios occur in
at most $4.5\%$ of component updates.  Hence the problem is not numerical
safeguarding; a
locally derived variance can decrease its own corrected scalar discrepancy
while increasing the coupled multi-center objective.  The joint-variance
block matches the legacy optimizer on clean data, but is slightly worse under
contamination and has a rare variance-collapse mode.  It is also 13--36\%
slower.  These results support retaining joint mean/variance updates within
each component as the primary optimizer.

\begin{table}[ht]
\centering
\scriptsize
\caption{Separated variance updates: mean (SD).  New optimizers use 20
datasets at variance ratio $1$, 30 at ratio $3$, and 20 in each contaminated
cell, with two or three model seeds.  Legacy rows are matched references from
the main experiments.}
\label{tab:variance_optimizers}
\resizebox{\linewidth}{!}{%
\begin{tabular}{@{}l l l c c c@{}}
\toprule
Data & Model & Optimizer & KL $\downarrow$ & Covariance $\downarrow$ & ARI $\uparrow$ \\
\midrule
\multirow{4}{*}{Clean, ratio $1$}
 & $3\Mone$ & Formula & .0872 (.0226) & .0748 (.0219) & .2998 (.0734) \\
 & $3\Mone$ & Joint variance & .2168 (.0688) & .3155 (.0909) & .2811 (.0639) \\
 & $3\Mone+8\Rtwo$ & Formula & .1161 (.0333) & .0813 (.0340) & \best{.3304 (.0640)} \\
 & $3\Mone+8\Rtwo$ & Joint variance & \best{.0855 (.0258)} & \best{.0457 (.0210)} & .3249 (.0618) \\
\addlinespace
\multirow{6}{*}{Clean, ratio $3$}
 & $\Mtwo$ & Formula & .4850 (.2014) & .3720 (.1221) & .2645 (.1342) \\
 & $\Mtwo$ & Joint variance & .1123 (.0310) & .0946 (.0453) & .5658 (.0625) \\
 & $\Mtwo$ & Legacy & .1112 (.0293) & .0925 (.0426) & .5689 (.0639) \\
 & $\Mone+4\Mtwo+4\Rtwo$ & Formula & .5708 (.1913) & .5289 (.3227) & .2374 (.1082) \\
 & $\Mone+4\Mtwo+4\Rtwo$ & Joint variance & \best{.1041 (.0279)} & \best{.0895 (.0408)} & .5765 (.0596) \\
 & $\Mone+4\Mtwo+4\Rtwo$ & Legacy & .1043 (.0277) & .0896 (.0389) & \best{.5788 (.0557)} \\
\addlinespace
\multirow{3}{*}{$5\%$ uniform}
 & $\Mtwo$ & Formula & .6426 (.1833) & .4654 (.2511) & .2182 (.1062) \\
 & $\Mtwo$ & Joint variance & .1484 (.0470) & .4097 (1.0234) & .5309 (.1001) \\
 & $\Mtwo$ & Legacy & \best{.1449 (.0384)} & \best{.1889 (.0880)} & \best{.5339 (.0800)} \\
\addlinespace
\multirow{3}{*}{$10\%$ uniform}
 & $\Mtwo$ & Formula & .8369 (.1253) & .5350 (.4467) & .1967 (.0746) \\
 & $\Mtwo$ & Joint variance & .3463 (.0696) & .5889 (.2051) & .3727 (.1190) \\
 & $\Mtwo$ & Legacy & \best{.3208 (.0656)} & \best{.5305 (.1883)} & \best{.4023 (.0978)} \\
\bottomrule
\end{tabular}}
\end{table}

The large covariance SD for the joint-variance block at $5\%$ is caused by a
single fit whose variance ratio exceeds $10^{12}$.  Excluding that dataset
reduces its mean covariance error to approximately $0.181$, close to the
legacy result; we retain the failure in the table because it is part of the
method's empirical stability.

\FloatBarrier
\subsection{Global heavy tails and asymmetric misspecification}
\label{sec:heavy_tails}

Uniform replacement creates sparse, spatially diffuse contamination.  The
Student-$t$ experiment instead makes every component heavy-tailed, so no
subset of observations is exactly Gaussian.  Table~\ref{tab:student_main}
shows that a small first-order contribution is now consistently useful.  At
$\nu=10$ and $5$, $\Mone+\Mtwo$ has significantly lower KL and higher ARI
than EM.  At $\nu=5$, KMeans has lower KL than the moment estimator but is
statistically tied in ARI.  At $\nu=3$, no likelihood-based Gaussian fit is a
good distributional approximation: KMeans has the lowest KL, the moment
method has much higher ARI than EM, but its KL and covariance error are worse.

\begin{table}[ht]
\centering
\scriptsize
\caption{Moment-matched Student-$t$ components: mean (SD) over 60 datasets
after averaging three model seeds.  The selected moment model uses weight one
on $\Mone$ for $\nu=10,5$ and weight $0.5$ for $\nu=3$.}
\label{tab:student_main}
\resizebox{\linewidth}{!}{%
\begin{tabular}{@{}c l c c c c c@{}}
\toprule
$\nu$ & Estimator & KL $\downarrow$ & Weight TV $\downarrow$ & Mean error $\downarrow$
 & Covariance $\downarrow$ & ARI $\uparrow$ \\
\midrule
\multirow{3}{*}{$10$}
 & $\Mone+\Mtwo$ & \best{.1451 (.0437)} & \best{.0654 (.0310)}
 & \best{.1588 (.0245)} & \best{.1468 (.0524)} & \best{.5180 (.0589)} \\
 & EM, $10^{-2}$ & .1595 (.0476) & .0951 (.0551) & .1618 (.0551)
 & .2382 (.0871) & .4859 (.0802) \\
 & KMeans & .2269 (.0532) & .1059 (.0359) & .2123 (.0342)
 & .2138 (.0535) & .4938 (.0669) \\
\addlinespace
\multirow{3}{*}{$5$}
 & $\Mone+\Mtwo$ & .3213 (.0770) & \best{.0720 (.0535)}
 & \best{.2483 (.5253)} & \best{.6724 (1.8032)} & .4554 (.0720) \\
 & EM, $10^{-2}$ & .4082 (.1003) & .1916 (.1005) & .5257 (.8856)
 & .7485 (.6956) & .3262 (.1116) \\
 & KMeans & \best{.2717 (.1110)} & .1220 (.0807) & .5327 (.9039)
 & 1.0237 (2.2632) & \best{.4670 (.1077)} \\
\addlinespace
\multirow{3}{*}{$3$}
 & $0.5\Mone+\Mtwo$ & 1.1560 (.4139) & \best{.1632 (.1302)}
 & \best{2.4679 (7.8186)} & 5.1856 (6.0362) & .4077 (.1255) \\
 & EM, $10^{-2}$ & 1.0063 (.2046) & .3455 (.1429) & 2.9965 (7.8140)
 & \best{1.7030 (.9426)} & .2024 (.1612) \\
 & KMeans & \best{.5654 (.2316)} & .2260 (.1643) & 3.0168 (7.8026)
 & 3.8213 (3.8819) & \best{.4280 (.2056)} \\
\bottomrule
\end{tabular}}
\end{table}
\FloatBarrier

Dataset-paired CIs qualify these comparisons.  Relative to EM, the moment
model changes KL/ARI by $-.0145$ ($[-.0238,-.0055]$) and $+.0321$
($[+.0181,+.0462]$) at $\nu=10$, and by $-.0869$
($[-.1057,-.0693]$) and $+.1292$ ($[+.1067,+.1531]$) at $\nu=5$.  At
$\nu=3$, it improves ARI by $+.2053$ ($[+.1761,+.2368]$) but worsens KL by
$+.1498$ ($[+.0576,+.2624]$).  Compared with KMeans, its ARI differences at
$\nu=5$ and $3$ are $-.0116$ and $-.0203$, and both intervals include zero.

The very large standard deviations at $\nu\leq5$ are caused by rare
degenerate fits, so means alone are misleading.  Table~\ref{tab:student_robust}
reports robust summaries and an explicit failure diagnostic.  At $\nu=5$ the
typical moment solution remains accurate: median mean error is $0.152$ and
median covariance error is $0.344$, far below their respective means.  At
$\nu=3$, however, the upper quartile is already affected; this is a genuine
stability limit rather than a handful of harmless outliers.

\begin{table}[ht]
\centering
\scriptsize
\caption{Heavy-tail stability of the selected $\Mone+\Mtwo$ model.  Errors
are median [IQR] over dataset-level averages.  A run is flagged when variance
ratio $>10^3$, minimum weight $<.01$, or mean error $>1$.}
\label{tab:student_robust}
\begin{tabular}{@{}c c c c c c@{}}
\toprule
$\nu$ & Mean error & Covariance error & Variance ratio
 & Failure: $\Mtwo$ & Failure: selected \\
\midrule
$10$ & .160 [.142,.176] & .145 [.114,.174] & 3.55 [3.23,3.82] & 0.0\% & 0.0\% \\
$5$  & .152 [.130,.166] & .344 [.285,.388] & 4.98 [4.47,5.85] & 8.3\% & 3.3\% \\
$3$  & .154 [.123,3.055] & .778 [.713,13.301] & 9.09 [6.26,$8.0\!\times\!10^{11}$]
 & 41.1\% & 36.7\% \\
\bottomrule
\end{tabular}
\end{table}
\FloatBarrier

The failures are strongly associated with initialization.  Conditional on a
failed KMeans initialization, pure $\Mtwo$ fails on $68.2\%$ of runs at
$\nu=5$ and $77.9\%$ at $\nu=3$; conditional on a stable initialization, no
failure is observed.  The largest squared Mahalanobis norm has Spearman
correlation $0.49$ and $0.73$ with failure at $\nu=5$ and $3$, respectively.
For 21 of the 60 datasets at $\nu=3$, all three moment seeds fail, so adding
random restarts alone cannot repair the severe regime.  These are statistical
failures even though every optimizer returns a nominally successful status.

The directional-Gamma experiment provides a complementary negative result.
Across all 12 asymmetric settings, regularized EM has the lowest true KL; in
the aggregate comparison its mean KL is $0.242$ (SD $0.036$ across settings),
versus $0.283$ (SD $0.066$) for the best moment configuration.  The paired
gap is positive in all 12 settings, with mean $0.040$ and SD $0.036$.
Localized moments therefore do not confer generic robustness to arbitrary
misspecification.  They are most effective when contamination is symmetric
within components or diffuse in space, and can be biased by coherent
directional skew.

\FloatBarrier
\clearpage
\subsection{Summary of empirical findings}
\label{sec:exp_summary}

The experiments support four conclusions.  First, on clean Gaussian data the
localized estimator is competitive in the homoscedastic and moderate-variance
regimes, but EM is more efficient when component scales differ strongly.
Second, directional second moments give a large and statistically stable
advantage under $2$--$10\%$ uniform contamination; the gain does not depend on
using outliers as test centers at $d\leq10$, but it does not survive the fixed-
feature-budget screen from $d=50$ onward.  Third, additional moment orders must be
treated as model-dependent regularizers: first moments help on clean and
globally heavy-tailed data but should be downweighted under diffuse
contamination, while zeroth moments and compensated radial moments are often
too sensitive to misspecification.  Finally, the default component-wise joint
mean/variance optimizer is not merely an implementation choice.  Formulaic
variance estimation is unreliable once scales differ, whereas a separated
joint variance block is broadly accurate but slower and slightly less stable.

The main limitation is severe global heavy tails.  In that regime the method
can preserve clustering substantially better than EM while producing rare
but extreme parameter estimates, and at $\nu=3$ KMeans is the better Gaussian
approximation by forward KL.  A robust initialization or explicit safeguards
against component-scale explosion are therefore necessary before the method
can be recommended as a general-purpose replacement for EM.  The dimension
screen adds a second limitation: the number and geometry of localized tests
must scale with the parameter dimension, not only with sample size.  All
configurations and per-run outputs are stored under
\texttt{experiments/configs/} and \texttt{results/}, with matching run names.

\end{document}
